
\documentclass{article}
\usepackage{colortbl}
\usepackage{makecell}
\usepackage{multirow}
\usepackage{supertabular}

\begin{document}

\newcounter{utterance}

\centering \large Interaction Transcript for game `ifeval', experiment `ifeval', episode 351 with Qwen3.5{-}4B.
\vspace{24pt}

{ \footnotesize  \setcounter{utterance}{1}
\setlength{\tabcolsep}{0pt}
\begin{supertabular}{c@{$\;$}|p{.15\linewidth}@{}p{.15\linewidth}p{.15\linewidth}p{.15\linewidth}p{.15\linewidth}p{.15\linewidth}}
    \# & \multicolumn{2}{c}{Player} && \multicolumn{2}{c}{Game Master} \\
    \hline

    \theutterance \stepcounter{utterance}  
    & & & \multicolumn{4}{p{0.6\linewidth}}{
        \cellcolor[rgb]{0.9,0.9,0.9}{
            \makecell[{{p{\linewidth}}}]{
                \texttt{\tiny{[P1$\langle$GM]}}
                \texttt{Write a parody of 'ars poetica'. Do not include the word 'parody' throughout your response.} \\
            }
        }
    }
     \\ \\

    \theutterance \stepcounter{utterance}  
    & \multicolumn{4}{p{0.6\linewidth}}{
        \cellcolor[rgb]{0.561,0.737,0.561}{
            \makecell[{{p{\linewidth}}}]{
                \texttt{\tiny{[P1$\rangle$GM]}}
                \texttt{**Ars Poetica: A Modern Guide to Digital Suffering**} \\
\\ 
\texttt{To write a poem that moves the reader,} \\
\texttt{You must first understand the reader's mood.} \\
\texttt{Do not begin with a grand, heroic mood,} \\
\texttt{But start with a notification sound.} \\
\\ 
\texttt{The Muse is dead; she lives in the cloud,} \\
\texttt{She waits for Wi{-}Fi to connect and load.} \\
\texttt{If your syntax is heavy and slow,} \\
\texttt{She will abandon you to the void.} \\
\\ 
\texttt{Do not describe a lion's roar,} \\
\texttt{But rather, a server that crashes hard.} \\
\texttt{Do not speak of ancient, golden lore,} \\
\texttt{But of a battery that dies at four.} \\
\\ 
\texttt{To make your verse truly sing,} \\
\texttt{You must include a broken link.} \\
\texttt{The reader should feel the sting,} \\
\texttt{Of a "404" error, stark and stark.} \\
\\ 
\texttt{Avoid the use of archaic words,} \\
\texttt{Like "thou" or "hath" or "ye."} \\
\texttt{Use "you" and "I" and "app,"} \\
\texttt{And let the algorithm decide.} \\
\\ 
\texttt{If you wish to evoke tears,} \\
\texttt{Do not describe a weeping willow tree.} \\
\texttt{Describe a friend who never replies,} \\
\texttt{Or a subscription that expires.} \\
\\ 
\texttt{The best poem is one that is short,} \\
\texttt{And ends with a "Like" button press.} \\
\texttt{The audience loves it, and it's worth it,} \\
\texttt{To see the engagement metrics rise.} \\
\\ 
\texttt{So write with care, and watch the screen,} \\
\texttt{For the digital muse is always near.} \\
\texttt{She feeds on data, and she's keen,} \\
\texttt{To turn your thoughts into a stream.} \\
            }
        }
    }
    & & \\ \\

    \theutterance \stepcounter{utterance}  
    & & & \multicolumn{2}{p{0.3\linewidth}}{
        \cellcolor[rgb]{0.9,0.9,0.9}{
            \makecell[{{p{\linewidth}}}]{
                \texttt{\tiny{[GM$|$GM]}}
                {'keywords:forbidden_words': {'forbidden_words': ['parody']}}
            }
        }
    }
    & & \\ \\

    \theutterance \stepcounter{utterance}  
    & & & \multicolumn{2}{p{0.3\linewidth}}{
        \cellcolor[rgb]{0.9,0.9,0.9}{
            \makecell[{{p{\linewidth}}}]{
                \texttt{\tiny{[GM$|$GM]}}
                \texttt{game\_result = WIN} \\
            }
        }
    }
    & & \\ \\

\end{supertabular}
}

\end{document}
